\subsection{Evaluation Protocol}
\label{sec:evaluation}

\paragraph{Clip-level event agreement.}
Valid listening reactions can occur at different times,
making strict alignment with a single reference overly
restrictive. We therefore compare event-category presence
over each clip, counting repeated occurrences only once.
We report macro-precision, macro-recall, and macro-F1
separately for visual and vocal categories with reference
positives, and report false-positive rates for the remaining
categories. These scores measure category agreement with
the observed response, complementing the assessment of
contextual appropriateness.

\paragraph{Audiovisual quality and behavior.}
Building on the OmniResponse evaluation
framework~\cite{omniresponse}, we measure video quality
using PSNR, SSIM, LPIPS, FID, and FVD.
Emo-Acc, Emo-Div, and $\Delta$Resp characterize expression
agreement, expression diversity, and listening-motion
differences, following AnyTalker~\cite{anytalker}.
Audio diagnostics measure activity-duration error,
missed vocal responses, excess vocal duration, and false
vocalization on silent clips. Where applicable, SyncNet's LSE-D and LSE-C
measure lip synchronization. Unsupported cases are
reported as N/A, together with valid evaluation coverage.

\paragraph{Human evaluation.}
Human evaluation targets contextual appropriateness,
naturalness and engagement, and audiovisual coordination
(\cref{sec:qualitative}). The assessment allows alternative
reasonable reactions and natural quiet listening.